\chapter{Restriction and intrinsic trace data}
\label{chap:trace-data}

Let $X$ be a connected compact complex manifold, let $Y\subset X$ be a connected compact complex submanifold of positive dimension, and fix a quasi-plurisubharmonic function $u\in\QPSH(X)$ whose generic Lelong number along $Y$ vanishes: $\nu(u,Y)=\inf_{y\in Y}\nu(u,y)=0$.

Even when all multiplier ideals of $u$ are trivial, its restriction to $Y$ can have logarithmic poles. To keep only information determined by the multiplier ideals of $u$ on $X$, we restrict the ideals themselves and take their asymptotic vanishing orders on modifications of $Y$. These orders form the divisorial trace data. Subadditivity of multiplier ideals and analytic approximation make the trace data additive under sums of potentials.

The next question is whether a potential on $Y$ realizes these numbers. On a curve, they form an atomic measure, whose Green potential has exactly the prescribed Lelong numbers. In higher dimension, restricted analytic envelopes converge to the largest potential satisfying the trace inequalities, but Lelong semicontinuity gives only a lower bound for its divisorial Lelong numbers. When $Y$ is K\"ahler, convergence in a sufficiently large closed background gives the reverse bound; pullback from a curve gives another exact realization, also for non-K\"ahler $Y$. The product examples at the end show that exact trace data still does not determine the Monge--Amp\`ere mass.

Choose a smooth real $(1,1)$-form $\theta$ with $u\in\PSH(X,\theta)$ and a Hermitian form $\omega$ on $X$. We use the analytic approximants $u_m$ of \cref{cor:nonclosed-analytic-approximation}: they have analytic type $(1/m,\mathcal I(mu))$ in the sense of \cref{sec:analytic-types} and satisfy $\theta+\ddc u_m\geq-\epsilon_m\omega$ with $\epsilon_m\downarrow0$. By \cref{cor:restricted-analytic-divisorial-bound}, their restrictions to $Y$ are quasi-plurisubharmonic, in particular not identically $-\infty$. As in \cref{chap:products}, $\mathrm{MA}^{\mathrm{ext}}$ denotes the extended non-pluripolar Monge--Amp\`ere measure, also written $\mathrm{MA}$ when its total mass is finite.

\section{Restricted ideals and trace data}\label{sec:restricted-ideals}\label{sec:trace-additivity}

The restriction of a multiplier ideal to $Y$ records the vanishing of its local generators along $Y$. Its vanishing order along a divisor on a modification of $Y$ can grow with the weight, even when the ideal is generically trivial on $Y$.

For $t>0$ put
\[
 \mathfrak a_t(u;Y)=\mathcal I(tu)\mathcal O_Y.
\]
Here $\mathcal I(tu)\mathcal O_Y$ denotes the image of $\mathcal I(tu)$ under the restriction map $\mathcal O_X\to\mathcal O_Y$. This is a coherent ideal sheaf on $Y$, and it is nonzero: since $\nu(u,Y)=0$, there is $y\in Y$ with $\nu(u,y)<1/t$, and Skoda's integrability theorem gives $\mathcal I(tu)_y=\mathcal O_{X,y}$.

In this chapter, a \emph{divisor over $Y$} is a prime divisor $E$ on a connected smooth compact modification $\pi\colon Z\to Y$. For a nonzero coherent ideal $\mathfrak a$ on $Y$, let $\operatorname{ord}_E(\mathfrak a)$ be the generic vanishing order along $E$ of the ideal generated by its pullback to $Z$; it is a nonnegative integer.

\Needspace{9\baselineskip}
\begin{theorem}\label{thm:intrinsic-trace-profile}\index{Trace!restricted ideal profile}
For every divisor $E$ over $Y$, the limit
\begin{equation}\label{eq:trace-profile}
 \tau_{Y,u}(E)
 :=\lim_{t\to\infty}\frac{\operatorname{ord}_E\mathfrak a_t(u;Y)}{t}
 =\sup_{t>0}\frac{\operatorname{ord}_E\mathfrak a_t(u;Y)}{t}
\end{equation}
exists and is finite, and it equals the limit along positive integers. These numbers depend only on the multiplier-ideal equivalence class of $u$ on $X$.
\end{theorem}
\begin{proof}
Fix the modification $\pi\colon Z\to Y$ carrying $E$, and write $b(t)=\operatorname{ord}_E\mathfrak a_t(u;Y)$ for $t>0$. For $s,t>0$, the subadditivity of multiplier ideals (the second inclusion in \eqref{eq:trace-local-ideal-inputs}), followed by restriction to $Y$, gives
\[
 \mathfrak a_{s+t}\subseteq\mathfrak a_s\mathfrak a_t,
 \qquad b(s+t)\geq b(s)+b(t).
\]
Moreover, $b(t)\geq0$. For fixed $s>0$ and $t=ks+r$ with $k$ a nonnegative integer and $0\leq r<s$, repeated superadditivity gives $b(t)\geq k b(s)$. Hence
\[
 \liminf_{t\to\infty}\frac{b(t)}t\geq\frac{b(s)}s.
\]
Taking the supremum over $s$ proves the existence of the limit and the supremum formula, with the value $+\infty$ allowed at this stage. In particular, the limit along positive integers is the same.

For every sufficiently large integer $m$, the restriction $u_m|_Y$ has analytic type $(1/m,\mathfrak a_m(u;Y))$ by \cref{prop:analytic-type-calculus}, and therefore
\[
 \frac{b(m)}m
 =\nu\bigl((u_m|_Y)\circ\pi,E\bigr).
\]
These numbers are uniformly bounded by \cref{cor:restricted-analytic-divisorial-bound}, so the limit along integers, and hence the supremum, is finite. Finally, the ideals $\mathfrak a_t(u;Y)$ are unchanged when $u$ is replaced by a potential with the same multiplier ideals, which proves the last assertion.
\end{proof}

We call the function $E\mapsto\tau_{Y,u}(E)$, defined on all divisors over $Y$, the \emph{divisorial trace data} of $u$ along $Y$.

\section{Addition of the intrinsic trace data}\label{sec:trace-additivity}

Return to a submanifold $Y$ of arbitrary dimension. Subadditivity of multiplier ideals gave one inequality for the trace of a sum. For the reverse inequality, first fix an analytic approximation to the sum, then compare it with finer approximations to the two summands using \cref{lem:analytic-minorant-product-ideals}.

\Needspace{8\baselineskip}
\begin{theorem}\label{thm:trace-data-additivity}
Let $u,v\in\QPSH(X)$ have zero generic Lelong number along the connected compact submanifold $Y$. Then, on every divisor over $Y$,
\[
 \tau_{Y,u+v}=\tau_{Y,u}+\tau_{Y,v}.
\]
\end{theorem}
\begin{proof}
The inequality $\geq$ was proved in \cref{prop:trace-profile-calculus}. Fix a sufficiently large integer $k$ and let $S_k$ be the global analytic approximant to $u+v$, of type $(1/k,\mathcal I(k(u+v)))$. By \cref{thm:demailly-hermitian-input}, $u+v\leq S_k$. Apply \cref{lem:analytic-minorant-product-ideals} with $s=S_k$ to the global analytic approximants $U_m,V_m$ of $u,v$, respectively, with types $(1/m,\mathcal I(mu))$ and $(1/m,\mathcal I(mv))$. Their restrictions and that of $S_k$ are genuine because all three original potentials have zero generic Lelong number along $Y$.

Restrict the resulting ordinary inequality to $Y$, pull back to any fixed smooth compact modification, and take the generic Lelong number along its divisor $E$. This gives
\[
 \frac{\operatorname{ord}_E\mathfrak a_m(u;Y)}m
 +\frac{\operatorname{ord}_E\mathfrak a_m(v;Y)}m
 \geq(1-\delta_m)
       \frac{\operatorname{ord}_E\mathfrak a_k(u+v;Y)}k.
\]
Let $m\to\infty$. Finiteness and the limit formula in \cref{thm:intrinsic-trace-profile} yield
\[
 \tau_{Y,u}(E)+\tau_{Y,v}(E)
 \geq\frac{\operatorname{ord}_E\mathfrak a_k(u+v;Y)}k.
\]
Now let $k\to\infty$. This proves the opposite inequality and hence equality.
\end{proof}

When $Y$ is a curve, the theorem gives
$D_{u+v}=D_u+D_v$. With one fixed nonpositive Green kernel the corresponding Green potentials add, by integration of positive measures. The atomic trace therefore satisfies
\[
 \operatorname{Tr}^{\mathrm{at}}_Y(u+v)
 \sim\operatorname{Tr}^{\mathrm{at}}_Y(u)
       +\operatorname{Tr}^{\mathrm{at}}_Y(v).
\]
Positive real homogeneity follows from homogeneity of the trace data and linearity of the Green operator.


\section{Comparison with ordinary restriction}

When the restriction $u|_Y$ is not identically $-\infty$, the restriction theorem for multiplier ideals bounds the trace data by its divisorial Lelong numbers. For analytic singularities, comparison with the approximants gives the reverse inequality.

\begin{proposition}\label{prop:trace-naive-comparison}\index{Trace!ordinary restriction}
If $v=u|_Y$ is not identically $-\infty$, then
\[
 \tau_{Y,u}(E)\leq\nu(v\circ\pi,E)
\]
for every divisor $E$ over $Y$. If $u$ has analytic singularities, equality holds.
\end{proposition}
\Needspace{6\baselineskip}
\begin{proof}
The restriction theorem, which is the first inclusion in \eqref{eq:trace-local-ideal-inputs}, and \cref{lem:bergman-divisorial-error}, applied on $Y$, give for every $t>0$
\[
 \operatorname{ord}_E\mathfrak a_t(u;Y)
 \leq\operatorname{ord}_E\mathcal I(tv)
 \leq t\nu(v\circ\pi,E).
\]
Divide by $t$ and use \cref{thm:intrinsic-trace-profile}.

If $u$ has analytic singularities, apply \cref{lem:analytic-minorant-product-ideals} to the pair $(u,0)$, with $s=u$, $U_m=u_m$ and $V_m=0$. It gives
\[
 u_m\leq(1-\delta_m)u+O(1),\qquad\delta_m\downarrow0.
\]
Restricting to $Y$, pulling back by $\pi$ and taking the generic Lelong number along $E$, we obtain, since $u_m$ has analytic type $(1/m,\mathcal I(mu))$,
\[
 \frac1m\operatorname{ord}_E\mathfrak a_m(u;Y)
 \geq(1-\delta_m)\nu((u|_Y)\circ\pi,E).
\]
Letting $m\to\infty$ proves the reverse inequality, hence equality. All the restrictions involved are quasi-plurisubharmonic, because $\nu(u,Y)=0$.
\end{proof}

The inequality can be strict, as the following local example shows. On a small bidisk in $\mathbb C^2$ with coordinates $(z,w)$, write
\[
 u(z,w)=\max\{-\log(-\log|z|^2),\ \log|w|^2\},
 \qquad Y=\{z=0\}.
\]
The bidisk is chosen with $|z|<\mathrm{e}^{-1}$, and the first function is extended by $-\infty$ along $z=0$. It is plurisubharmonic, since $s\mapsto-\log(-s)$ is increasing and convex for $s<0$, and hence so is the maximum. For every $t>0$,
\[
 \mathrm{e}^{-tu}\leq(-\log|z|^2)^t,
\]
which is locally integrable. Thus all multiplier ideals of $u$ are trivial and the trace data vanishes, whereas $u|_Y=\log|w|^2$ has Lelong number one at the origin. This example is taken from \cite[Example~8.1.2]{SGPT}.

\Needspace{12\baselineskip}
\section{Realization on a compact curve}\label{sec:curve-trace-realization}

We return to the compact setting, and until \cref{sec:trace-envelope} we assume that $Y$ is a curve. Since every modification of a compact smooth curve is an isomorphism, the divisors over $Y$ are its points, and the trace data prescribes an atomic curvature measure. Stokes' theorem bounds its total mass by the degree of the background on $Y$, and filling the remaining mass with a smooth measure gives a positive current whose potential realizes the data.

Set
\[
 D_u=\sum_{p\in Y}\tau_{Y,u}(p)\,\delta_p,
\]
where the sum is understood as the supremum of its finite subsums.

\Needspace{17\baselineskip}
\begin{theorem}\label{thm:curve-trace-realization}\index{Trace!realization on a curve}
Put $M=\int_Y\theta|_Y$. Then
\begin{equation}\label{eq:curve-trace-budget}
 0\leq D_u(Y)\leq M.
\end{equation}
The finite atomic measure $D_u$ determines a canonical ordinary singularity type of quasi-plurisubharmonic functions on $Y$: for every smooth positive form $\eta$ on $Y$ with $\int_Y\eta=1$, this type has a representative $v\in\PSH(Y,\theta|_Y)$ with
\begin{equation}\label{eq:curve-trace-current}
 \theta|_Y+\ddc v=D_u+(M-D_u(Y))\eta.
\end{equation}
The extended Monge--Amp\`ere measure of $v$ is $(M-D_u(Y))\eta$. The ordinary type does not depend on $\eta$, on the choice of the smooth form $\theta$ with $u\in\PSH(X,\theta)$, or on the choice of $u$ within its multiplier-ideal class.
\end{theorem}
\begin{proof}
For each positive integer $m$, let
\[
 D_m=\frac1m\sum_{p\in Y}
       \operatorname{ord}_p\mathfrak a_m(u;Y)\,\delta_p.
\]
A nonzero coherent ideal on a compact smooth curve has finitely many zeros, so this sum is finite. Since $u_m|_Y$ has analytic type $(1/m,\mathfrak a_m(u;Y))$, the measure $D_m$ is the atomic part of the positive current
$(\theta+\epsilon_m\omega)|_Y+\ddc(u_m|_Y)$, and
Stokes' theorem on the curve gives
\[
 D_m(Y)\leq M+\epsilon_m\int_Y\omega|_Y.
\]
For any finite subset of $Y$, passing to the limit in the sum of the corresponding coefficients by \cref{thm:intrinsic-trace-profile}, and then taking the supremum over finite subsets, gives \eqref{eq:curve-trace-budget}. In particular, the positive coefficients have at most countable support and define a finite measure. Write $d=D_u(Y)$.

Choose a Green kernel $G$ for $\ddc$ on the curve, normalized so that
\[
 \ddc_xG(x,p)=\delta_p-\eta,\qquad G(x,p)\leq0.
\]
Such a kernel exists: the Green kernel of the scalar Laplacian on the compact Riemann surface has, in a local coordinate $z$, the singularity $\log|z(x)-z(p)|^2$, is uniformly bounded above, and satisfies $\sup_p\int_Y|G(x,p)|\eta(x)<\infty$; subtracting a constant gives the sign condition. By Fubini's theorem,
\[
 w(x)=\int_YG(x,p)\,\mathrm{d}D_u(p)
\]
is integrable, and it is quasi-subharmonic with
$\ddc w=D_u-d\,\eta$. Its Lelong number at every point $p$ is the mass of $D_u$ at $p$, namely $\tau_{Y,u}(p)$.

Since $M\eta-\theta|_Y$ is smooth with integral zero, the Poisson equation gives a smooth $s$ with $\ddc s=M\eta-\theta|_Y$. Then $v=w+s$ satisfies \eqref{eq:curve-trace-current}, and the right-hand side is positive because $M\geq d$. At every atom of $D_u$, the positive Lelong number forces $v=-\infty$. The extended measure of $v$ is the restriction of the current in \eqref{eq:curve-trace-current} to the set where $v>-\infty$; this set contains no atom of $D_u$, and its complement is polar and therefore not charged by $\eta$. Hence the extended measure is $(M-d)\eta$.

If another smooth form $\eta$ or another smooth background $\theta$ is used, the two constructed representatives have the same atomic curvature, and all other terms in their equations are smooth. Their difference therefore has smooth $\ddc$, so it is smooth by elliptic regularity, and in particular bounded. This proves independence of these choices at the level of ordinary types. Independence of the choice of $u$ within its multiplier-ideal class holds because the numbers defining $D_u$ depend only on that class.
\end{proof}

We denote this ordinary type by $\operatorname{Tr}^{\mathrm{at}}_Y(u)$ and call it the \emph{atomic trace} of $u$ on the curve $Y$; the same symbol also denotes any of its representatives. All potentials on a curve with these Lelong numbers have the same multiplier ideals, but they need not have the same ordinary type; the construction selects the type whose curvature, apart from the atoms, is a smooth measure. Additivity of the trace data (\cref{thm:trace-data-additivity}) and linearity of the Green operator give
\[
 \operatorname{Tr}^{\mathrm{at}}_Y(u+v)
 \sim\operatorname{Tr}^{\mathrm{at}}_Y(u)+\operatorname{Tr}^{\mathrm{at}}_Y(v),
 \qquad
 \operatorname{Tr}^{\mathrm{at}}_Y(au)\sim a\operatorname{Tr}^{\mathrm{at}}_Y(u)
 \quad(a>0).
\]
Here $v\in\QPSH(X)$ is any function that also has zero generic Lelong number along $Y$.

Keep $D_m$, $G$ and $w$ from the proof of \cref{thm:curve-trace-realization}, and put $w_m=\int_YG(\cdot,p)\,\mathrm{d}D_m(p)$.

\begin{proposition}\label{prop:curve-trace-approximation}
The measures $D_m$ satisfy $D_m\leq D_u$ and converge to $D_u$ in total variation, and they increase along dyadic indices. For a fixed nonpositive Green kernel, the potentials $w_{2^j}$ decrease to $w$, and $w_m\to w$ in $L^1(Y)$. Each $w_m$ has the ordinary singularity type of $u_m|_Y$.

If the restriction $u|_Y$ is not identically $-\infty$, it is at least as singular as the atomic trace:
\[
 u|_Y\preceq \operatorname{Tr}^{\mathrm{at}}_Y(u).
\]
If $u$ has analytic singularities, the two are ordinarily equivalent.
\end{proposition}
\begin{proof}
The supremum formula in \eqref{eq:trace-profile} gives $D_m\leq D_u$ coefficient by coefficient. Since the coefficients converge and are dominated by those of the finite measure $D_u$, dominated convergence gives $D_m(Y)\to D_u(Y)$, and hence convergence in total variation. Superadditivity with $s=t=m$ gives $D_{2m}\geq D_m$, and since $G\leq0$ the potentials $w_{2^j}$ decrease. Monotone convergence for $-G$ identifies their pointwise limit with $w$, including at poles. The uniform $L^1$ bound for the kernel also gives
\[
 \|w_m-w\|_{L^1(\eta)}
 \leq C\|D_m-D_u\|_{\mathrm{TV}}\longrightarrow0.
\]
Each $D_m$ is finitely supported, and locally both $w_m$ and $u_m|_Y$ are the sum of the logarithmic singularities prescribed by $D_m$ and a bounded function. Thus their ordinary types agree.

If $a=u|_Y$ is not identically $-\infty$, \cref{prop:trace-naive-comparison} gives $\nu(a,p)\geq\tau_{Y,u}(p)$ at every point. Since $\ddc w=D_u-D_u(Y)\eta$, \cref{lem:divisorial-potential-domination} implies $a\preceq w$. If $u$ has analytic singularities, \cref{prop:trace-naive-comparison} identifies the finitely many logarithmic coefficients of $u|_Y$ with those of $D_u$, and the bounded local remainders give ordinary equivalence.
\end{proof}

\section{The Poisson equation in the atomic trace type}

Write $d=D_u(Y)$, and let $w$ be the Green representative constructed in the proof of \cref{thm:curve-trace-realization}. For every $\theta|_Y$-plurisubharmonic potential in the ordinary type of $w$, the part of its curvature measure carried by polar sets is $D_u$, and its extended measure is the remaining part, of total mass $M-d$. Prescribing the extended measure therefore amounts to solving a Poisson equation for the bounded difference from $w$.

\begin{theorem}\label{thm:curve-atomic-lp-equation}
Let $\mathrm{d}V$ be a smooth positive volume form on $Y$, and let $0\leq f\in L^p(\mathrm{d}V)$, $p>1$, with $\int_Yf\,\mathrm{d}V>0$.
There is a solution $v$ in the atomic trace type of
\[
 \mathrm{MA}^{\mathrm{ext}}_{\theta|_Y}(v)=c f\,\mathrm{d}V,
 \qquad c>0,
\]
if and only if $M>d$. In that case
\[
 c=\frac{M-d}{\int_Yf\,\mathrm{d}V},
\]
and $v$ is unique up to an additive constant. If $M=d$, there is instead a potential in this type with extended measure zero, unique up to an additive constant.
\end{theorem}
\begin{proof}
Assume that $M>d$, take the value of $c$ given above, and solve
\[
 \ddc s=d\,\eta+c f\,\mathrm{d}V-\theta|_Y.
\]
The right-hand side has integral zero, so the elliptic $L^p$ theory gives a solution $s\in W^{2,p}(Y)$, unique up to constants. In real dimension two, the Sobolev embedding for $p>1$ makes $s$ continuous and bounded. The potential $v=w+s$ has the atomic trace type and satisfies
\[
 \theta|_Y+\ddc v=D_u+c f\,\mathrm{d}V\geq0.
\]
Its extended measure is $c f\,\mathrm{d}V$, since this absolutely continuous measure does not charge the polar set of poles, while the atoms lie in it. For $M=d$, the same construction with $c=0$ has a smooth right-hand side and gives the asserted potential with zero extended measure.

To prove necessity and uniqueness, let $a$ be any $\theta|_Y$-plurisubharmonic potential in the ordinary type of $w$, and put $\mu=\theta|_Y+\ddc a$. By \cref{cor:polar-order-maximum}, ordinary equivalence with $w$ makes the polar part of $\mu$ equal to $D_u$. Consequently
\[
 \mathrm{MA}^{\mathrm{ext}}_{\theta|_Y}(a)=\mu-D_u,
 \qquad
 \int_Y\mathrm{MA}^{\mathrm{ext}}_{\theta|_Y}(a)=M-d.
\]
Integrating the equation therefore forces the stated value of $c$, so a positive $c$ requires $M>d$. If two potentials solve the equation with the same density and constant, their curvature currents both equal $D_u+c f\,\mathrm{d}V$; their difference has zero $\ddc$ and is therefore constant on the compact connected curve. The same argument proves uniqueness in the case of zero extended measure.
\end{proof}

\section{Comparison with the source metric on the curve}

Every smooth positive form on a compact Riemann surface is K\"ahler. We may therefore use the source book's metric on singularities on $Y$, while the ambient manifold $X$ remains arbitrary Hermitian. This gives a quantitative comparison for the trace just constructed.

\begin{proposition}\label{prop:curve-trace-source-distance}
Choose a smooth positive form $\eta$ of integral one on $Y$ and choose $A>D_u(Y)$ sufficiently large that every restricted approximant $u_m|_Y$ is $A\eta$-plurisubharmonic. Use a fixed nonpositive Green kernel to define $w,w_m$ as above. In the source book's normalization of $d_{S,A\eta}$,
\begin{equation}\label{eq:curve-trace-source-distance}
 d_{S,A\eta}(w,u_m|_Y)
 =d_{S,A\eta}(w,w_m)
 =\frac12\bigl(D_u(Y)-D_m(Y)\bigr)
 \longrightarrow0.
\end{equation}
In particular the atomic trace has $\mathcal I$-good singularities in the source sense on the curve. If the ambient manifold is K\"ahler, it represents the $P$-equivalence class of the source book's trace.
\end{proposition}
\begin{proof}
The restricted curvature bounds of \cref{cor:nonclosed-analytic-approximation} are dominated by a single $A\eta$, after increasing $A$. Since $D_m\leq D_u$ and $G\leq0$, we have $w\leq w_m$. Also
\[
 A\eta+\ddc w=D_u+(A-D_u(Y))\eta,
 \qquad
 A\eta+\ddc w_m=D_m+(A-D_m(Y))\eta.
\]
Their extended masses are $A-D_u(Y)$ and $A-D_m(Y)$, both positive. We now invoke only the following explicit closed-background identity from \cite[chapter on comparison of singularities, ordered-distance formula]{SGPT}: for $a\preceq_P b$ in a fixed closed background $\Theta$ on a compact K\"ahler manifold of dimension $n$,
\[
 d_{S,\Theta}(a,b)
 =\frac1{n+1}\sum_{j=0}^n
 \left(\int\Theta_b^j\wedge\Theta_{V_\Theta}^{n-j}
       -\int\Theta_a^j\wedge\Theta_{V_\Theta}^{n-j}\right),
\]
where the products are non-pluripolar. In dimension one the $j=0$ terms cancel. Substitution of the two masses proves the second equality in \eqref{eq:curve-trace-source-distance}. Ordinary equivalence of $u_m|_Y$ and $w_m$ gives distance zero between them in the source theory, proving the first equality. Total-variation convergence of $D_m$ proves the limit.

Each $w_m$ has finitely many logarithmic poles with rational coefficients, hence analytic singularities. The positive-mass characterization recalled in \cref{chap:hermitian-igood}, now applied on the K\"ahler curve $Y$, shows that their $d_S$ limit $w$ is $\mathcal I$-good.

Finally suppose $X$ is K\"ahler. The source book supplies a quasi-equisingular approximation with analytic type
$(2^{-j},\mathcal I(2^ju))$. Its restriction to $Y$ has the ordinary type of $w_{2^j}$. The same distance calculation therefore proves that these restrictions converge to $w$ in the source metric. This is precisely its definition of a trace representative. Uniqueness of that trace modulo $P$-equivalence gives the final assertion.
\end{proof}

\Needspace{5\baselineskip}
The metric used here lives on the curve. The proof does not construct a metric on singularities for an arbitrary higher-dimensional Hermitian background. It identifies exactly the part of the source trace theory that the atomic construction already recovers.


\section{The trace envelope}\label{sec:trace-envelope}

Return to a submanifold $Y$ of arbitrary positive dimension. The vanishing curvature error places an $L^1$ limit of normalized restricted approximants in the original background $\theta|_Y$. Its divisorial Lelong numbers dominate the trace data. Taking the largest potential with these inequalities gives a canonical envelope, which is also the limit of the ordinary envelopes of the restricted approximants. Enlarging the background then supplies the positive mass needed for prescribed-type equations.

\begin{proposition}\label{prop:trace-profile-majorant}
There is a potential $v\in\PSH(Y,\theta|_Y)$ such that
\[
 \nu(v\circ\pi,E)\geq\tau_{Y,u}(E)
 \quad\hbox{for every divisor }E\hbox{ over }Y.
\]
In particular, $v\in\PSH(Y,\Omega_0)$ for every Hermitian form $\Omega_0\geq\theta|_Y$.
\end{proposition}
\begin{proof}
The restrictions $u_m|_Y$ have a common smooth curvature lower bound. Subtract their suprema and extract an $L^1$ limit $v$ by normalized compactness. Their current inequalities pass to the limit as $\epsilon_m\to0$, giving $\theta|_Y+\ddc v\geq0$.

\Needspace{8\baselineskip}
For each fixed compact modification, \cref{lem:modification-lone-pullback} gives convergence of the pullbacks along that same subsequence. \Cref{lem:lelong-lone-semicontinuity} gives, at every point $z$ of $E$,
\[
 \nu(v\circ\pi,z)
 \geq\limsup_m\nu((u_m|_Y)\circ\pi,z)
 \geq\lim_m\frac{\operatorname{ord}_E\mathfrak a_m(u;Y)}m
 =\tau_{Y,u}(E).
\]
Taking the generic value proves the assertion. No simultaneous extraction over divisors is required: the pullback lemma applies to each modification after the subsequence on $Y$ has been fixed.
\end{proof}

\Needspace{10\baselineskip}
For a smooth real $(1,1)$-form $\Omega$ on $Y$ for which the following family is nonempty, define
\begin{equation}\label{eq:trace-profile-envelope}
 \begin{split}
 \mathcal T_\Omega(u;Y):=\Bigl(\sup\{a\in\PSH(Y,\Omega):\ &a\leq0,\\
 &\nu(a\circ\pi,E)\geq\tau_{Y,u}(E)\
 \text{ for every }E\}\Bigr)^*.
 \end{split}
\end{equation}
The choice of $\Omega$ is part of this notation. The definition refers only to the intrinsic trace data, and therefore to the ambient multiplier-ideal class of $u$. The preceding proposition makes it available for $\Omega=\theta|_Y$ and for every larger background.

\subsection{Limits of restricted analytic envelopes}

The restricted ideal $\mathfrak a_m(u;Y)$ has a log resolution on a compact modification of $Y$. Its finitely many resolved coefficients are bounded by the trace data. Thus every trace competitor is more singular than the restricted analytic approximant. This comparison identifies the envelope limit even when an exact realization of the trace numbers is unavailable.

\begin{theorem}\label{thm:restricted-analytic-envelope-limit}
Let $\Omega$ be a smooth real $(1,1)$-form on $Y$, let $h$ be Hermitian, and choose $\delta_m\geq0$ tending to zero such that, for all sufficiently large $m$,
\[
 a_m:=u_m|_Y\in\PSH(Y,\Omega+\delta_m h).
\]
Then $T=\mathcal T_\Omega(u;Y)$ is genuine, and
\begin{equation}\label{eq:restricted-analytic-envelope-limit}
 q_m:=P_{\Omega+\delta_m h}[a_m]
     =P_{\Omega+\delta_m h}[a_m]_{\mathcal I}
       \longrightarrow T\quad\text{in }L^1(Y).
\end{equation}
If $\delta_{2^j}$ decreases, then $q_{2^j}\searrow T$ pointwise. The limit depends only on $\Omega$ and the ambient multiplier-ideal class of $u$. Moreover,
\[
 P_\Omega[T]_{\mathcal I}=P_\Omega[T]=T.
\]
\end{theorem}
\begin{proof}
The argument of \cref{prop:trace-profile-majorant}, applied to the assumed curvature bounds, gives a normalized $\Omega$-plurisubharmonic limit $v$ of a subsequence of the $a_m$. Its divisorial numbers dominate $\tau_{Y,u}$. Hence the family defining $T$ is nonempty, and its supremum lies between $v$ and zero. It is stable under finite maxima, since Lelong numbers of a maximum are the minima of the two Lelong numbers, also after pullback. Choquet extraction and finite maxima give an increasing sequence of competitors with canonical limit $T$. Pullback convergence and Lelong semicontinuity show that $T$ retains every defining inequality. Thus $T$ is itself a competitor and is the largest one.

Fix $m$ and resolve $\mathfrak a_m(u;Y)$ on $\pi:Z\to Y$. For each prime component $E$ of the resolved ideal, \eqref{eq:trace-profile} gives
\[
 \nu(T\circ\pi,E)\geq\tau_{Y,u}(E)
       \geq\frac1m\operatorname{ord}_E\mathfrak a_m(u;Y)
       =\nu(a_m\circ\pi,E).
\]
The finite-divisor criterion in \cref{lem:analytic-ideal-dominance} gives $T\preceq a_m$. Since $T\leq0$ and $\delta_m\geq0$, it is a competitor for $q_m$. Also \cref{prop:analytic-envelope-type} identifies the two envelopes in \eqref{eq:restricted-analytic-envelope-limit} and gives $q_m\sim a_m$. Therefore
\[
 T\leq q_m\leq0,\qquad
 \nu(q_m\circ\pi,E)=\frac1m\operatorname{ord}_E\mathfrak a_m(u;Y)
\]
for every divisor over $Y$.

All $q_m$ have supremum zero and a common smooth curvature lower bound. Consider any $L^1$ cluster limit $q$. It belongs to $\PSH(Y,\Omega)$, lies between $T$ and zero, and has
\[
 \nu(q\circ\pi,E)\geq\limsup_m\nu(q_m\circ\pi,E)
       =\tau_{Y,u}(E)
\]
along the convergent subsequence. Here the pullback and semicontinuity argument of \cref{prop:trace-profile-majorant} applies to each fixed compact modification. Thus $q$ is a competitor for $T$, which proves $q=T$. Normalized compactness now gives convergence of the full sequence in $L^1$.

Subadditivity gives $\mathfrak a_{2m}(u;Y)\subseteq\mathfrak a_m(u;Y)^2$, hence $a_{2m}\preceq a_m$. When the dyadic backgrounds decrease, this implies $q_{2^{j+1}}\leq q_{2^j}$. The lower bound by $T$ makes the pointwise limit genuine, and its $L^1$ identification gives equality with $T$ everywhere. Finally the defining family for $T$ depends only on $\Omega$ and the ideals of $u$, proving the asserted independence.

Ideal projection preserves all divisorial Lelong numbers, so $P_\Omega[T]_{\mathcal I}$ is a trace competitor. Maximality and the envelope inequality give $P_\Omega[T]_{\mathcal I}=T$. The ordinary fixed-point assertion follows from \cref{cor:ordinary-preserves-i}.
\end{proof}

Taking $\Omega=\theta|_Y$, $h=\omega|_Y$ and $\delta_m=\epsilon_m$ gives the limit in the original restricted background. Along dyadic indices the errors can be enlarged to a decreasing sequence with limit zero. A fixed sufficiently large Hermitian background permits $\delta_m=0$.

Inclusion of the competitor families also orders the trace envelopes: increasing the prescribed data decreases $\mathcal T_\Omega$, whereas increasing the background increases it whenever both envelopes are defined.

\subsection{Positive mass and prescribed equations}

\begin{theorem}\label{thm:trace-profile-model-envelope}
Choose $\Omega_0,v$ as in \cref{prop:trace-profile-majorant} and set $\Omega=2\Omega_0$. Then $T=\mathcal T_\Omega(u;Y)$ is genuine and satisfies all the divisorial inequalities in \eqref{eq:trace-profile-envelope}. Moreover,
\[
 P_\Omega[T]_{\mathcal I}=T,\qquad P_\Omega[T]=T,
 \qquad P_{\Omega/2}(T)\not\equiv-\infty.
\]
Writing $r=\dim Y$, its extended measure satisfies
\[
 0<\int_Y\mathrm{MA}^{\mathrm{ext}}_\Omega(T)<+\infty,
 \qquad
 \mathrm{MA}^{\mathrm{ext}}_\Omega(T)
 \leq\mathbf1_{\{T=0\}}\Omega^r.
\]
\end{theorem}
\begin{proof}
Normalize $v\leq0$. It is a trace competitor in background $\Omega$, so the defining supremum lies between $v$ and zero. The finite-maxima and ideal-projection arguments in the proof of \cref{thm:restricted-analytic-envelope-limit} apply to this family. They give the genuine envelope $T$, all its defining inequalities and both fixed-point identities.

The original $\Omega_0$-plurisubharmonic $v$ is below $T$, so it supplies a competitor for $P_{\Omega/2}(T)$. The fixed-point contact estimate \cref{prop:fixed-point-finite}, applied on $Y$, gives the displayed measure bound and finiteness. The smaller-background individual positivity in \cref{lem:smaller-background-mass-lower} gives strict positivity.
\end{proof}

\Needspace{13\baselineskip}
\begin{corollary}\label{cor:trace-envelope-prescribed-equation}
For the chosen $\Omega$ and $T$, let $\mathrm{d}V$ be a smooth positive volume form on $Y$ and let $0\leq f\in L^p(\mathrm{d}V)$, $p>1$, with $\int_Y f\,\mathrm{d}V>0$. For every $\lambda>0$ there is a unique $a\sim T$ in $\PSH(Y,\Omega)$ with
\[
 \mathrm{MA}^{\mathrm{ext}}_\Omega(a)=\mathrm{e}^{\lambda a}f\,\mathrm{d}V.
\]
There is also a positive scalar $c$ and $a\sim T$ with $\mathrm{MA}^{\mathrm{ext}}_\Omega(a)=cf\,\mathrm{d}V$; the scalar is unique.
\end{corollary}
\begin{proof}
The preceding theorem checks both parts of the model condition \eqref{eq:als-model-condition} on the compact Hermitian manifold $Y$. Apply the prescribed ordinary-type theorem \cref{thm:als-prescribed-lp}.
\end{proof}

On a curve this model construction can be compared more precisely with the atomic trace. Take $\Omega=A\eta$ as in \cref{prop:curve-trace-source-distance}, and let $w$ be its Green representative. By \cref{lem:divisorial-potential-domination}, the defining pointwise trace inequalities for $a\in\PSH(Y,A\eta)$ are equivalent to $a\preceq w$. Thus the defining competitor family is exactly the ordinary singularity-envelope family, and
\[
 \mathcal T_{A\eta}(u;Y)=P_{A\eta}[w]=P_{A\eta}[w]_{\mathcal I}.
\]
The last equality follows from the $\mathcal I$-goodness proved in \cref{prop:curve-trace-source-distance}. In particular the model envelope has exactly the prescribed curve Lelong numbers.


\section{Exact traces over a curve}
\label{sec:curve-fibration-traces}

An exact representative on a curve also gives higher-dimensional examples by pullback. Multiplier ideals commute with pullback by a submersion, so the trace data on a modification is given by the trace data on the curve, multiplied by the multiplicities of the pulled-back fibers. The pulled-back Green potential has precisely these divisorial Lelong numbers, and its domination property identifies the ordinary type of the trace envelope.

In this section, $p:X\to M$ is a surjective holomorphic submersion of connected compact complex manifolds with connected fibers. Let $C\subset M$ be a connected compact smooth curve, put $Y=p^{-1}(C)$, and write $q=p|_Y:Y\to C$; then $Y$ is a connected compact submanifold of $X$. Take $U\in\QPSH(M)$ with $\nu(U,C)=0$, and set $u=U\circ p$. The generic Lelong number of $u$ along $Y$ is zero, since in local product coordinates it equals that of $U$ along $C$. Write $\tau_{C,U}$ for the trace data of $U$ along $C$, defined as in \cref{sec:restricted-ideals} with $(M,C,U)$ in place of $(X,Y,u)$, and let
\[
 D=\sum_{b\in C}\tau_{C,U}(b)\,\delta_b,
 \qquad \delta=D(C)<+\infty.
\]
Choose a smooth positive form $\eta$ of integral one on $C$ and a nonpositive Green kernel for $\eta$. The construction in the proof of \cref{thm:curve-trace-realization} gives $w\in\QPSH(C)$, not identically $-\infty$, with
\[
 \ddc w=D-\delta\eta.
\]
\Needspace{9\baselineskip}
For $b\in C$, put $Y_b=q^{-1}(b)$, and let $[Y_b]$ be its current of integration. Set $v=w\circ q$. Pullback of distributions by a submersion is defined by duality with integration along the fibers; it commutes with differentiation and, in the present situation, with positive sums, so
\begin{equation}\label{eq:curve-fiber-divisorial-current}
 \ddc v=\sum_{b\in C}\tau_{C,U}(b)[Y_b]-\delta q^*\eta.
\end{equation}
Each fiber is a smooth connected prime divisor, and the sum is a positive closed current of finite mass, as also follows locally from the positivity of $\ddc(w\circ q)+\delta q^*\eta$.

\Needspace{24\baselineskip}
\begin{theorem}\label{thm:curve-fibration-exact-trace}\index{Trace!curve fibration}
The potential $v$ realizes the entire divisorial trace data of $u$ along $Y$:
\[
 \nu(v\circ\pi,E)=\tau_{Y,u}(E)
\]
for every prime divisor $E$ on every smooth compact modification $\pi:Z\to Y$. More explicitly, write $h=q\circ\pi$. If $h(E)=C$, both sides are zero. If $h(E)=\{b\}$, choose a local coordinate $t_b$ on $C$ vanishing at $b$; then both sides equal
\[
 \operatorname{ord}_E(t_b\circ h)\,\tau_{C,U}(b).
\]

Choose any Hermitian form $\Omega_0$ with $v\in\PSH(Y,\Omega_0)$, and set $\Omega=2\Omega_0$. Then
\begin{equation}\label{eq:curve-fibration-trace-envelope}
 \mathcal T_\Omega(u;Y)=P_\Omega[v]=P_\Omega[v]_{\mathcal I}\sim v.
\end{equation}
The envelope has positive finite extended mass, and the conclusions of \cref{cor:trace-envelope-prescribed-equation} on equations in its ordinary type hold. Its ordinary type does not depend on the choice of the Green kernel, of the form $\eta$, or of the sufficiently large background $\Omega$.
\end{theorem}
\begin{proof}
\Cref{lem:submersion-multiplier-pullback} and restriction to $Y$ give
\[
 \mathfrak a_t(u;Y)=\mathfrak a_t(U;C)\mathcal O_Y.
\]
The nonzero coherent ideal $\mathfrak a_t(U;C)$ on the smooth curve is generated near $b$ by $t_b^{k_t(b)}$, where $k_t(b)=\operatorname{ord}_b\mathfrak a_t(U;C)$. For a divisor $E$ with $h(E)=\{b\}$, the order along $E$ of the extended ideal is therefore $k_t(b)\operatorname{ord}_E(t_b\circ h)$. For a divisor dominating $C$, the generic order is zero, since the ideal on the curve has only finitely many zeros. Dividing by $t$ and letting $t\to\infty$ gives exactly the stated formula for $\tau_{Y,u}(E)$.

For the Lelong numbers of $v\circ\pi$, suppose first that $h(E)=\{b\}$. At a general point of $E$, a local equation $z_1=0$ of $E$ satisfies $t_b\circ h=z_1^k g$ with $g$ holomorphic and nonvanishing, and $k=\operatorname{ord}_E(t_b\circ h)$. Replacing $z_1$ by $z_1g^{1/k}$ for a local holomorphic $k$th root, we may assume $t_b\circ h=z_1^k$. The Lelong number of the one-variable function $w(z_1^k)$ at $0$ is $k\nu(w,b)=k\tau_{C,U}(b)$, as follows from the definition through circular means, and this is also the generic Lelong number of $v\circ\pi$ along $E$ in these product coordinates. If $h(E)=C$, the regular locus of $E$ contains a point mapped to a point where $w$ is finite, because the nonconstant holomorphic map from that locus to $C$ has open image; the pullback is finite there, so its generic Lelong number along $E$ is zero. This proves exact realization on every compact model.

By \cref{lem:divisorial-potential-domination} and \eqref{eq:curve-fiber-divisorial-current}, a quasi-plurisubharmonic function $a$ on $Y$ satisfies $a\preceq v$ if and only if $\nu(a,Y_b)\geq\tau_{C,U}(b)$ for every $b$ with $\tau_{C,U}(b)>0$. These fiber inequalities are among the trace inequalities, and conversely domination by $v$ implies all trace inequalities on modifications, by the exact formula just proved. Thus the trace competitors for $\Omega$ are exactly the competitors for the ordinary envelope of $v$, and $\mathcal T_\Omega(u;Y)=P_\Omega[v]$.

For $\psi=P_\Omega[v]_{\mathcal I}$, the Lelong numbers along the fibers are those of $v$, since $\psi$ and $v$ have the same multiplier ideals, and the same domination lemma gives $\psi\preceq v$. Hence
\[
 v-\sup_Yv\leq P_\Omega[v]\leq\psi\preceq v.
\]
Since $\psi\leq0$, it is a competitor for the ordinary envelope of $v$, so $\psi\leq P_\Omega[v]$. This proves \eqref{eq:curve-fibration-trace-envelope}.

The potential $v-\sup_Yv$ is a competitor for $P_{\Omega/2}(P_\Omega[v])$. The contact estimate for envelope fixed points, \cref{prop:fixed-point-finite}, gives finite mass, and \cref{lem:smaller-background-mass-lower} gives strict positivity. The two envelope properties give the model condition \eqref{eq:als-model-condition}, so \cref{thm:als-prescribed-lp} gives the stated equations.

Two Green constructions on $C$ differ by a smooth function, by the Poisson argument in the proof of \cref{thm:curve-trace-realization}, so their pullbacks have the same ordinary type. Every background in the statement gives an envelope ordinarily equivalent to this pullback. This proves the asserted independence, at the level of ordinary types rather than of normalized potentials.
\end{proof}

The same construction respects sums and multiplication by positive real numbers at the level of ordinary types. Indeed, the trace measures on the curve add by \cref{thm:trace-data-additivity}, their Green representatives add for a fixed kernel, and pullback preserves sums. Each envelope in a sufficiently large background represents the corresponding ordinary type.


\section{Product realizations and their measures}
\label{sec:trace-product-measures}

For a sum of potentials pulled back from the two factors of a product with a curve, dimension reasons force every term of the Monge--Amp\`ere expansion to contain exactly one factor from the curve. The resulting product formula computes the measure of the pulled-back Green representative and solves equations with separated densities. Adding a potential from the other factor with trivial multiplier ideals can keep every trace coefficient while changing finite mass to infinite mass.

In this section, $C$ is a compact connected curve with a smooth positive form $\eta$ of integral one, and $(Z,\omega_Z)$ is a compact connected Hermitian manifold of dimension $k\geq1$. On $Y=C\times Z$, let $q:Y\to C$ and $r:Y\to Z$ be the projections. Choose $A>0$ and set
\[
 \Omega=q^*(A\eta)+r^*\omega_Z.
\]
For positive Borel measures $\mu$ on $C$ and $\nu$ on $Z$, $\mu\otimes\nu$ denotes their product measure. The extended measures below are $\sigma$-finite by \cref{thm:arbitrary-extended-ma}, so this notation is meaningful also when a total mass is infinite.

\begin{proposition}\label{prop:curve-product-extended-ma}\index{Trace!product measure}
For $w\in\PSH(C,A\eta)$ and $h\in\PSH(Z,\omega_Z)$, the potential $a=w\circ q+h\circ r$ is $\Omega$-plurisubharmonic and
\begin{equation}\label{eq:curve-product-extended-ma}
 \mathrm{MA}^{\mathrm{ext}}_\Omega(a)
 =(k+1)\,
       \mathrm{MA}^{\mathrm{ext}}_{A\eta}(w)
                \otimes\mathrm{MA}^{\mathrm{ext}}_{\omega_Z}(h).
\end{equation}
\end{proposition}
\begin{proof}
The curvature current of $a$ is the sum of the pullbacks of the two positive currents $A\eta+\ddc w$ and $\omega_Z+\ddc h$, which proves positivity.

Suppose first that both potentials are bounded. Expanding the bounded product by multilinearity, only the terms with one factor from $C$ and $k$ factors from $Z$ survive: two factors from $C$ vanish for dimension reasons, and so do $k+1$ factors from $Z$. The binomial coefficient is $k+1$, which gives the formula for bounded potentials once the exterior product $q^*(A\eta+\ddc w)\wedge r^*(\omega_Z+\ddc h)^k$ is identified with the product measure $(A\eta+\ddc w)\otimes(\omega_Z+\ddc h)^k$. To justify this identification at the present regularity, add local plurisubharmonic corrections on each factor and regularize by decreasing convolution on smaller product charts. For the smooth approximants the identity is the exterior product of forms. Local Bedford--Taylor convergence, together with the finite expansion in the smooth background, gives weak convergence of the factor measures, with uniformly bounded variation on compact subcharts. The product measures then converge weakly: one first tests against products of continuous functions on the two factors and then uses uniform approximation by finite sums of such products. This proves the identity in the limit.

For general potentials, subtract constants so that $w,h\leq0$, and set $w_j=\max(w,-j)$ and $h_j=\max(h,-j)$. On
\[
 R_j=\{w>-j\}\times\{h>-j\}
\]
the sum agrees with the bounded potential $w_j\circ q+h_j\circ r$. This set is plurifine open and Borel, and it lies in $\{a>-2j\}$. Locality of the products of bounded truncations therefore identifies the restriction to $R_j$ of the extended measure of $a$ with that of the bounded sum. By the bounded case, this is $k+1$ times the product of the restrictions of the truncated measures, which are exactly the restrictions of the two extended measures of the factors.

The sets $R_j$ increase to the product of the complements of the two pole sets, and both sides of \eqref{eq:curve-product-extended-ma} give zero mass to the remaining set: for the left-hand side it is the pole set of $a$, and for the right-hand side this follows because each factor measure charges no pluripolar set and the product measures are $\sigma$-finite. Monotone convergence of positive measures proves the formula, even when the total mass is infinite.
\end{proof}

Suppose now that $D$ is a finite atomic measure on $C$, for instance the trace measure of \cref{sec:curve-fibration-traces}, put $\delta=D(C)$, and let $w$ be the Green potential with $\ddc w=D-\delta\eta$. For $A>\delta$, the potential $v=w\circ q$ satisfies
\begin{equation}\label{eq:product-trace-smooth-measure}
 \mathrm{MA}^{\mathrm{ext}}_\Omega(v)
       =(k+1)(A-\delta)\,\eta\otimes\omega_Z^k.
\end{equation}
Thus its measure is smooth and positive, although $v$ may have infinitely many singular fibers. This identity concerns the extended measure, whose definition discards the divisorial part of the curvature current.

The product formula also solves equations with separated densities in this ordinary type. Fix a smooth positive volume form $\mathrm{d}V_Z$ on $Z$ and $p>1$, and let $0\leq f\in L^p(C,\eta)$ and $0\leq g\in L^p(Z,\mathrm{d}V_Z)$ have positive integrals. By \cref{thm:curve-atomic-lp-equation}, applied on $C$ in the background $A\eta$, there is $w_f\sim w$ with
\[
 \mathrm{MA}^{\mathrm{ext}}_{A\eta}(w_f)=c_C f\eta,
 \qquad c_C=\frac{A-\delta}{\int_Cf\eta}.
\]
\Needspace{6\baselineskip}
By \cref{thm:kn-continuous}, there are a bounded $h_g\in\PSH(Z,\omega_Z)$ and $c_Z>0$ with
\[
 \mathrm{MA}_{\omega_Z}(h_g)=c_Zg\,\mathrm{d}V_Z.
\]
The sum $a=w_f\circ q+h_g\circ r$ is ordinarily equivalent to $v$, and
\[
 \mathrm{MA}^{\mathrm{ext}}_\Omega(a)
       =(k+1)c_Cc_Z\,(f\eta)\otimes(g\,\mathrm{d}V_Z).
\]
The product density belongs to $L^p$ by Fubini's theorem. This is an explicit construction for separated densities; existence for general $L^p$ densities, in a sufficiently large background, follows from \cref{cor:trace-envelope-prescribed-equation}.

\begin{proposition}\label{prop:exact-trace-infinite-representatives}
Suppose that $Z$ does not have the bounded mass property, that is, $v_+(\omega_Z)=+\infty$ in the notation of \cref{def:bmp}, and that $A>\delta$. Then there is $a\in\PSH(Y,\Omega)$ such that
\[
 a\sim_{\mathcal I}v,\qquad
 \int_Y\mathrm{MA}^{\mathrm{ext}}_\Omega(v)<+\infty,
 \qquad
 \int_Y\mathrm{MA}^{\mathrm{ext}}_\Omega(a)=+\infty.
\]
Both potentials have the same divisorial Lelong numbers on every modification; when $D$ is the trace measure of \cref{sec:curve-fibration-traces}, these are the trace data computed in \cref{thm:curve-fibration-exact-trace}. The product manifold $Y$ does not have the bounded mass property.
\end{proposition}
\begin{proof}
Choose the nonpositive potential $H$ of \cref{thm:infinite-mass-witness} on $Z$; it has infinite extended mass, and $\mathrm{e}^{-tH}$ is integrable for every $t>0$. Put $a=w\circ q+H\circ r$. The product formula and $A-\delta>0$ show that the mass of $a$ is infinite, while \eqref{eq:product-trace-smooth-measure} gives the finite positive mass of $v$.

We prove equality of the multiplier ideals locally, for every $t>0$. Since $H\leq0$, integrability against $\mathrm{e}^{-ta}$ implies integrability against $\mathrm{e}^{-tv}$, so $\mathcal I(ta)\subseteq\mathcal I(tv)$. Conversely, \cref{lem:submersion-multiplier-pullback} gives $\mathcal I(tv)=\mathcal I(tw)\mathcal O_Y$. Local generators pulled back from $C$ are square integrable against $\mathrm{e}^{-tw}$, so by Fubini's theorem they are also square integrable against $\mathrm{e}^{-tw}\mathrm{e}^{-tH}=\mathrm{e}^{-ta}$; multiplication by bounded holomorphic functions and finite sums preserve this local integrability. Thus $\mathcal I(tv)\subseteq\mathcal I(ta)$, which proves equality at every positive weight. \Cref{thm:local-i-inputs} then gives equality of all divisorial Lelong numbers after modification.

Finally, since $a$ has infinite mass, \cref{prop:bmp-finite} shows that $Y$ does not have the bounded mass property.
\end{proof}

For a concrete ambient realization, take a compact complex manifold $M$ containing the curve $C$, a potential $U\in\QPSH(M)$ with $\nu(U,C)=0$, and $X=M\times Z$. Then $Y=C\times Z$ and $u=U\circ\operatorname{pr}_M$, where $\operatorname{pr}_M:X\to M$ is the projection, satisfy the hypotheses of \cref{thm:curve-fibration-exact-trace} with $p=\operatorname{pr}_M$. The factor $Z$ can be the Hopf threefold of \cref{thm:hopf-unbounded}. Small logarithmic poles of $U$ at finitely many points of $C$ already give nonzero trace data, and the construction on the curve also allows any countable trace data arising from the multiplier ideals of $U$.

The mass formula of \cref{sec:resolved-trace-masses} uses more information: the restriction itself in a closed background, with a finite divisorial part and a bounded remainder on a K\"ahler model. Its residual class is nef and computes the mass of that restriction. Divisorial trace data alone does not provide such a remainder, as the preceding example shows.

